\section{模型假设}

结合题目附录 2、附录 3 的统一计算口径，本文做出如下假设，每条均在后文使用：

\begin{enumerate}[leftmargin=2em]
  \item \textbf{航段几何假设}：任意两任务节点间以两节点水平直线为运输航段，计划巡航海拔取该航段沿途 $30\,$m DEM 像元最高地面高程再加 $50\,$m 净空。该假设与附录 2 一致，用于统一水平距离、爬升高度与下降高度的计算。
  \item \textbf{三段式飞行假设}：每个航段依次经历爬升、巡航、下降三阶段，$O01$ 作业高度取其地面海拔，服务区作业高度取地面海拔以上 $30\,$m，连续多点访问时每次投送后从 $30\,$m 重新爬升。下降能耗效率取 $0$，即不单独计入下降附加能耗。
  \item \textbf{等效航程线性假设}：机型 $g$ 携带载荷 $q$ 时的等效航程 $L_g(q)$ 在空载标准航程与满载标准航程之间随载荷线性插值，用以将载荷折算为能耗，符合附录 2 规则。
  \item \textbf{能源与充电假设}：电池/能源组件初始 SOC 均为 $100\%$，采用两阶段等效充电模型（$0\%\!\sim\!90\%$ 占等效充电时间 $65\%$，$90\%\!\sim\!100\%$ 占 $35\%$），同一资源任务占用与充电时段不重叠、不同资源可并行；任务结束剩余 SOC 不得低于返航电量下限。
  \item \textbf{货箱不可拆分假设}：每个货箱为不可拆分整体，只能被安排一次，且同一货箱不跨架次拆分；组批仅受载质量与装载体积两类容量约束。
  \item \textbf{通信判定假设}：通信端点位置由三维坐标确定，链路按双向链路预算判定，取两方向最大允许损耗的较小值；地形遮挡按 $30\,$m DEM 视线连线是否被地面阻挡附加 $L_{obs}$ 损耗，中继不做多跳转发。
  \item \textbf{独立执行假设（问题四）}：任务分区后各组独立执行，运输无人机、共享电池、中继无人机与中继能源组件均不得跨组调配，各组分别核算资源需求。
\end{enumerate}
